\documentclass[11pt]{article}
\usepackage[margin=.9in]{geometry}
\usepackage{amsmath,amssymb,amsthm}
\usepackage[T1]{fontenc}
\usepackage{lmodern,microtype,hyperref}
\newtheorem{theorem}{Theorem}
\newtheorem{lemma}[theorem]{Lemma}
\newtheorem{proposition}[theorem]{Proposition}
\newtheorem{corollary}[theorem]{Corollary}
\theoremstyle{remark}\newtheorem{remark}[theorem]{Remark}
\newcommand{\norm}[1]{\left\lVert#1\right\rVert}
\DeclareMathOperator{\diag}{diag}\DeclareMathOperator{\ran}{ran}
\newcommand{\R}{\mathbb R}\newcommand{\C}{\mathbb C}
\begin{document}
% Notebook 14 body fragment, no preamble. Unique label prefix: nc-.
% Requires amsmath, amssymb, amsthm and a one-argument \norm macro.

\section{The exact change of reference in the averaging differential}
\label{nc-section}

The derivative of the corrected average with fixed reference one
is not the same coordinate operator as the source average $Q(V)$
based at a nonflat reference $V$. We now compute their relation,
including the endpoint term. The block ratio $L$ is fixed, all
logarithms use their common identity branch, and all Hilbert norms
use physical volume. This calculation is internally audited at the
level of the pinned deposited transcriptions; no fresh verification
of the original CMP98 typesetting is claimed.

\subsection{Two references and an exact group identity}
\label{nc-group}

Write $M$ for the ordinary group average. Raw I:46965--46970
defines
\[
 \Phi^{\mathfrak g}(V,B;y,x)=\frac1i\log[
   (e^{iB}V)\langle y,x\rangle V\langle y,x\rangle^{-1}],
 \qquad
 \Phi^{\flat}(V,B;c)=\frac1i\log[M(e^{iB}V)(c)M(V)(c)^{-1}].
\]
For $U$ near $V$ set
\begin{equation}
 g_y^V(U)=\exp\left(i\sum_{x\in B(y)}L^{-d}
       \frac1i\log[U\langle y,x\rangle V\langle y,x\rangle^{-1}]
                   \right).
 \label{nc-endpoint-mean}
\end{equation}
The bridge I:47076--47087, expressly incorporating CMP98 (120),
identifies the corrected group increment as
\begin{equation}
 C_V(U)(c)=g_{c_-}^V(U)^{-1}M(U)(c)g_{c_+}^V(U)M(V)(c)^{-1}.
 \label{nc-moving-corrected}
\end{equation}
It equals one at $U=V$. With the logarithmic/linear split
I:24344--24349, the same bridge gives
\begin{equation}
 C_V(e^{iB}V)=e^{i\mathcal Q(V,B)},\qquad
 \mathcal Q(V,B)=LQ(V)B+O(B^2).
 \label{nc-source-linearization}
\end{equation}
This is the deposited identification of the middle coordinate;
the earlier transcription does not itself reproduce CMP98 (120).
We use the explicit bridge, rather than guessing that omitted input
formula.

Set $a_y(U)=g_y^1(U)$. The corrected map with fixed identity
reference is
\begin{equation}
 C_1(U)(c)=a_{c_-}(U)^{-1}M(U)(c)a_{c_+}(U).
 \label{nc-fixed-corrected}
\end{equation}
For fixed $V$ define
$k_y(U;V)=g_y^V(U)^{-1}a_y(U)a_y(V)^{-1}$. Direct substitution
and cancellation, preserving the factor order, give
\begin{equation}
 \begin{split}
 C_1(U)(c)C_1(V)(c)^{-1}
 =\operatorname{Ad}_{a_{c_-}(V)^{-1}}\bigl[
 k_{c_-}(U;V)^{-1}C_V(U)(c)
       \operatorname{Ad}_{M(V)(c)}k_{c_+}(U;V)\bigr].
 \end{split}
 \label{nc-group-identity}
\end{equation}
In particular no multiplicative cocycle for the endpoint mean has
been assumed.

For $U(s)=e^{isB}V$ let
\[
 \eta_y^1(V;B)=\frac1i\left.\frac d{ds}\right|_0
       a_y(U(s))a_y(V)^{-1},\qquad
 \eta_y^V(B)=\frac1i\left.\frac d{ds}\right|_0g_y^V(U(s)),
 \qquad Z_VB(y)=\eta_y^1(V;B)-\eta_y^V(B).
\]
Define the \emph{unscaled} endpoint difference by
\begin{equation}
 d_Wz(c)=z(c_-)-\operatorname{Ad}_{W(c)}z(c_+).
 \label{nc-endpoint-difference}
\end{equation}
This explicit definition fixes the sign independently of other
covariant-derivative conventions. All middle factors in
\eqref{nc-group-identity} equal one at $s=0$, so differentiating
gives
\begin{equation}
 \frac1i\left.\frac d{ds}\right|_0
 C_1(e^{isB}V)(c)C_1(V)(c)^{-1}
 =\operatorname{Ad}_{a_{c_-}(V)^{-1}}
       [LQ(V)B(c)-d_{M(V)}(Z_VB)(c)].
 \label{nc-angular-derivative}
\end{equation}

As a separate check, I:46977 defines
$LQ^{\flat}(V)B=\Phi_1^{\flat}(V;B)$. Differentiating the
general bridge directly at $B=0$ gives
\begin{equation}
 LQ(V)B=LQ^{\flat}(V)B-d_{M(V)}\overline r(V;B),
 \qquad
 \overline r(V;B)(y)=\sum_xL^{-d}\Phi_1^{\mathfrak g}(V;B;y,x).
 \label{nc-rooted-versus-corrected}
\end{equation}
Thus the later rooted operator $Q^{\flat}$ and the corrected source
operator $Q$ are distinct. Substitution into
\eqref{nc-angular-derivative} reproduces the direct derivative of
\eqref{nc-fixed-corrected}, confirming its endpoint sign.

\subsection{An explicit formula for the endpoint correction}
\label{nc-jacobians}

For a small complex Lie-algebra element $X$ put
\begin{equation}
 \mathcal J_XK=\int_0^1\operatorname{Ad}_{e^{itX}}K\,dt,
 \qquad
 \frac1i\left.\frac d{ds}\right|_0e^{i(X+sK)}e^{-iX}
           =\mathcal J_XK.
 \label{nc-exponential-jacobian}
\end{equation}
The integral power series and a Neumann series show that
$\mathcal J_X$ and its inverse are analytic near zero, and
$\norm{\mathcal J_X-I}+\norm{\mathcal J_X^{-1}-I}\le C\norm{X}$.
This does not require unitarity.

Write
\[
 V\langle y,x\rangle=e^{iY_x},\qquad
 \overline Y_y=\sum_xL^{-d}Y_x,\qquad
 r_x(V;B)=\Phi_1^{\mathfrak g}(V;B;y,x).
\]
The last quantity includes the actual contour orientations,
transports and repetitions. Applying
\eqref{nc-exponential-jacobian} to each contour logarithm and then
to the exponential of their mean proves
\begin{equation}
 Z_VB(y)=\mathcal J_{\overline Y_y}
       \sum_xL^{-d}\mathcal J_{Y_x}^{-1}r_x(V;B)
          -\sum_xL^{-d}r_x(V;B).
 \label{nc-Z-formula}
\end{equation}
In particular $Z_1=0$. The leading background term is
\[
 Z_VB=\frac i2\left(
   [\overline Y,\overline r]-\sum_xL^{-d}[Y_x,r_x]\right)
   +O\bigl((\max_x|Y_x|)^2\max_x|r_x|\bigr).
\]
There is no general identity forcing this difference of means of
commutators to vanish. The endpoint correction in
\eqref{nc-angular-derivative} must therefore be retained.

\subsection{The physical derivative conversion and its bound}
\label{nc-physical}

At fine spacing $h$ write $V=e^{ihA}$ and define the fixed-reference
physical map by
$C_1(e^{ihA})(c)=e^{iLh\mathcal F_h(A)(c)}$. For an ordinary
potential variation $K$, the fine angular left increment is
$h\mathcal J_{hA}K$ and the coarse left increment is
$Lh\mathcal J_{Lh\mathcal F_h(A)}D\mathcal F_h(A)K$.
Substituting these into \eqref{nc-angular-derivative} gives
\begin{equation}
 \begin{split}
 \mathcal J_{Lh\mathcal F_h(A)}D\mathcal F_h(A)K
 =\operatorname{Ad}_{a_-(V)^{-1}}\bigl[
 Q_h(V)\mathcal J_{hA}K
  -L^{-1}d_{M(V)}Z_V(\mathcal J_{hA}K)\bigr].
 \end{split}
 \label{nc-physical-unrearranged}
\end{equation}
Consequently the exact operator conversion is
\begin{equation}
 \boxed{\displaystyle
 Q_h(V)=\operatorname{Ad}_{a_-(V)}
   \mathcal J_{Lh\mathcal F_h(A)}D\mathcal F_h(A)
        \mathcal J_{hA}^{-1}+L^{-1}d_{M(V)}Z_V.}
 \label{nc-physical-conversion}
\end{equation}
Here $Q_h(V)$ is the normalized source coefficient placed between
the physical spaces at $h$ and $Lh$. Neither $d_{M(V)}$ nor $Z_V$
contains an inverse mesh factor. The argument of $Z_V$ before
division is $h\mathcal J_{hA}K$; its linearity cancels $h$ against
the coarse $Lh$ and leaves exactly $1/L$.

\begin{proposition}[Actual one-step source comparison]
\label{nc-one-step-proposition}
For fixed $L$, if $h\norm{A}_\infty$ lies in a sufficiently small
common complex neighborhood, then
\begin{equation}
 \norm{Q_h(e^{ihA})-Q_h(1)}_{p,h\to p,Lh}
     \le C_{L,p}h\norm{A}_\infty,
 \qquad p=2,\infty.
 \label{nc-one-step-bound}
\end{equation}
The constants are independent of $h$, refinement and compatible
periodic volume. The same estimate holds for absolute coefficient
row sums and physical-adjoint row sums, after changing the fixed
constant.
\end{proposition}

\begin{proof}
Each one-step contour has bounded length at fixed $L$. Products of
possibly nonunitary complex transports therefore give
$\max_x|Y_x|+|\overline Y|\le C_Lh\norm{A}_\infty$ and
$\norm{r_x(V;\cdot)}_{\infty\to\mathrm{color}}\le C_L$.
Equation \eqref{nc-Z-formula} implies
$\norm{Z_V}_{\infty\to\infty}\le C_Lh\norm{A}_\infty$.

For the physical Hilbert bound expand $Z_V$ into its finitely many
one-step bond coefficients. Each is analytic, vanishes at $V=1$,
and is bounded by $C_Lh\norm{A}_\infty$. Rows and columns have
fixed finite incidence; the physical adjoint supplies the fixed
volume ratio $L^d$. The Schur bound thus gives the same order
from fine physical $L^2$ to coarse-site physical $L^2$, with no
macroscopic path-trace estimate. The endpoint map $d_{M(V)}$ has
two entries per row, fixed column incidence and bounded complex
adjoints, so it is uniformly bounded in both norms.

The other coordinate factors in \eqref{nc-physical-conversion}
differ from the identity by at most $C_Lh\norm{A}_\infty$.
Combine this with the proved one-step estimate
$\norm{D\mathcal F_h(A)-Q_h(1)}_{p\to p}
\le D_ph\norm{A}_\infty$ and the flat contractions to obtain
\eqref{nc-one-step-bound}. Equivalently, the entire corrected
source coefficient in \eqref{nc-source-linearization} is a finite
local analytic function of the link logarithms $hA$, equal to
$Q_h(1)$ at zero. Its coefficientwise Taylor bound and the same
finite incidence prove the row and physical-adjoint row assertions
directly. Formula \eqref{nc-physical-conversion} records the
coordinate factors hidden by that shorter proof. Periodic
multiplicity counting preserves all constants.
\end{proof}

For a multilevel source comparison this estimate must be evaluated
at the actual ordinary reference backgrounds $M^t(U)$, whose
amplitudes are controlled in the following argument. It does not
license replacing them by the corrected fixed-reference iterates.
It also proves no nonflat inverse or minimizing-field theorem by
itself.

\end{document}
